%------------------------------------------------------------------------------%

\usepackage{apresentacao}
\usepackage{tabto}

\graphicspath{{./figs/}}

%------------------------------------------------------------------------------%
\newcommand{\mytitle}   {Convergência das Séries de Taylor}
\title   {\mytitle}
\author  {\large Luis Alberto D'Afonseca}
\date    {}

%------------------------------------------------------------------------------%
\begin{document}

%------------------------------------------------------------------------------%
\begin{frame}
    \titlepage
\end{frame}

%------------------------------------------------------------------------------%
\section{Convergência da Série de Taylor}

%------------------------------------------------------------------------------%
\begin{frame}{Polinômios e Séries de Taylor}
  Se $f$ é $k$-derivável, em uma vizinhança de $\;x=a$, temos o \emph{Polinômio de Taylor}
  \[
    P_k(x) = \sum_{n=0}^k \frac{\;f^{(n)\;}(a)}{n!}\,(x-a)^n
  \]

  \pause
  \vspace{\baselineskip}
  Se $f$ é infinitamente derivável temos a \emph{Série de Taylor}
  \[
    \tilde{f}(x) = \sum_{n=0}^\infty \frac{\;f^{(n)\;}(a)}{n!}\,(x-a)^n
  \]
\end{frame}

%------------------------------------------------------------------------------%
\begin{frame}{Questões}
  \begin{itemize}[<+->]
    \item Quão próximo o polinômio $\;P_k(x)\;$ está da função $\;f(x)\;$?
    \vspace{\baselineskip}
    \item A série converge, isso é, $\;\tilde{f}(x)\;$ existe?
    \vspace{\baselineskip}
    \item Se $\;\tilde{f}(x)\;$ existe, qual sua relação com $\;f(x)\;$?
  \end{itemize}
\end{frame}

%------------------------------------------------------------------------------%
\begin{frame}{Teorema de Taylor}
  Sejam
  \begin{itemize}
  \item $\;f:\mathcal{I}\rightarrow \R\;$ uma função com \emph{$n+1$ derivadas contínuas} em $\mathcal{I}$
  \item $a$, $x$ $\in\mathcal{I}$
  \end{itemize}

  \pause
  \vspace{\baselineskip}
  Então existe {\color{red}$\;c\;$} \emph{entre $\;a\;$ e $\;x\;$} \pause tal que
  \[
    f(x)=P_n(x)+R_n(x)\
    \qquad \text{onde} \qquad
    R_n(x)=\frac{\;f^{(n+1)}({\color{red}c})\;}{(n+1)!}(x-a)^{n+1}
  \]

  \vspace{.5\baselineskip}
  $R_n(x)\;$ é o \emph{Resto de Ordem $n$}
\end{frame}

%------------------------------------------------------------------------------%
\begin{frame}{Convergência da Série de Taylor}
Se \emph{\(\;R_n(x)\to 0\;\)} quando \emph{\(\;n\to\infty\;\)}
para todo \(\;x\in\mathcal{I}\;\)

\vspace{\baselineskip}
A série de Taylor gerada por $f$ converge para $f$ em \(\;\mathcal{I}\;\)
\[
  f(x) = \sum_{n=0}^\infty \frac{\;f^{(n)\;}(a)}{n!}\,(x-a)^n
\]
\end{frame}

%------------------------------------------------------------------------------%
\begin{frame}{Estimativa de Erro}
Se existir \(\;M\;\) tal que  \emph{\(\;\abs{\,f^{(n+1)}(t)} \leq M \;\)}
para todo $t$ entre $x$ e $a$, inclusive

\pause
\vspace{\baselineskip}
\[
  \abs{R_n(x)}
  = \abs{\frac{\;f^{(n+1)}({\color{red}c})\;}{(n+1)!}(x-a)^{n+1}} \pause
  \leq M\, \frac{\;\abs{x-a}^{n+1}\;}{(n+1)!}
\]
\end{frame}

%------------------------------------------------------------------------------%
\section{Exemplo}

%------------------------------------------------------------------------------%
\begin{frame}{Exemplo}
   Obter uma aproximação para \(e^x\) no intervalo \([-1, 1]\)

   \vspace{0.5\baselineskip}
   \pause
   Série de Taylor da exponencial
   \[
     f(x) = \sum_{n=0}^\infty \frac{x^n}{n!}
   \]

   \pause
   Erro após a soma de $n$ termos
   \[
     R_n(x)=\frac{\;f^{(n+1)}(c)\;}{(n+1)!}x^{n+1}
   \]
   Para algum $c$ entre zero e $x$
\end{frame}

%------------------------------------------------------------------------------%
\begin{frame}{Exemplo}
  Não temos como saber o valor de $c$\pause, mas sabemos que
  \[
    f^{(n+1)}(c) = e^c \pause \leq e^1 \pause < 3
  \]
  \pause
  Portanto
  \[
    \abs{R_n(x)} < \frac{\;3\abs{x^{n+1}}\;}{(n+1)!} \pause
                 < \frac{\;3\;}{(n+1)!}
  \]
\end{frame}

%------------------------------------------------------------------------------%
\begin{frame}{\mytitle}
  \tableofcontents
\end{frame}

%------------------------------------------------------------------------------%
\end{document}
%------------------------------------------------------------------------------%
